\section{\secmathind}
\label{sec_math_ind}

In this section, we introduce mathematical induction, a powerful proof technique used throughout mathematics and theoretical computer science to prove a statement $P(n)$ that depends on an integer $n$. For example, it is used in computer science to verify loop invariants. We practice using mathematical induction to prove some of our conjectures about explicit rules for recursively-defined sequences from \Cref{sec_recursively_defined_sequences}.   

%%%%%%%%%%%%%%%%%%%%%%%%

\newcommand{\subpffmathind}{Proof Format: Mathematical Induction}
\subsection{\subpffmathind}
\label{sub_pff_math_induction}

Mathematical induction is used to prove a statement $P(n)$ for all values of the integer $n \ge s$ for some nonnegative integer $s$.  That is, to prove the infinite list of statements $P(s)$, $P(s+1)$, $P(s+2)$, \ldots. A proof by mathematical induction takes an inventive approach.

\begin{defn}[Proof by mathematical induction]
\label{defn_math_ind}
  A proof by \vocab{mathematical induction} that a statement $P(n)$ is true for all integers $n \ge s$, where $s$ is a nonnegative integer, consists of two parts.
      \begin{apart}
          \item First, we prove the \vocab{base case}: $P(s)$.
          \item  Next, we prove the \vocab{inductive step}: for any integer $n$, if $n > s$ and $P(n-1)$, then $P(n)$.
          \item To prove the inductive step, we assume $P(n-1)$ for some integer $n > s$ and prove $P(n)$ using a direct proof of the conditional (\Cref{pff_if_then_direct}).  The assumption of $P(n-1)$ is the \vocab{inductive hypothesis}.
      \end{apart}
\end{defn}

Here is the official proof format.

\begin{pff}[Mathematical induction]
\label{pff_math_ind}
    Prove $P(n)$ for $n \ge s$.
    \begin{proof}
        Induct on $n$. First, when $n=s$ we have (explain why $P(s)$ is true).

        Next, assume $P(n-1)$ for some integer $n > s$. That is, (restate $P(n-1)$ more clearly, if needed). Then (explain why $P(n)$ is true).
    \end{proof}
\end{pff}

The phrase ``induct on $n$'' alerts the reader that we are using proof by mathematical induction. The key words in the proof are ``First'' (indicates the base case), ``Assume'' (indicates the inductive hypothesis), ``for Some'' (indicates the existential quantifier), and ``Then'' (indicates the end of the inductive step).  A mnemonic is that a proof by mathematical induction is FAST.

Why does this proof format work? The result of the inductive step is $P(n-1) \to P(n)$ whenever $n > s$.  That is, if we know that $P(n-1)$ is true for some integer $n >s$, then we can conclude that $P(n)$ is true. Since we have proved $P(s)$ in the base case, we can apply the result of the inductive step with $n=s+1$ to get $P(s) \to P(s+1)$, from which it follows that $P(s+1)$.  Officially we are using the logical inference \emph{modus ponens} from \Cref{sec_inference}. Then, since we know $P(s+1)$, we can apply the result of the inductive step with $n=s+2$ to get $P(s+1) \to P(s+2)$ from which it follows that $P(s+2)$. And so on.

Try reading a proof by mathematical induction.

\begin{act}[Divisibility proof by mathematical induction]
\label{act_div_proof_math_ind}
~
    \begin{actpart}
        \item Copy the following proof.

        Use mathematical induction to prove that $4 \mid (5^n-1)$ for all integers $n \ge 1$
            \begin{proof}
                First, when $n=1$ we have $5^1-1=4$ so $4 \mid (5^1-1)$.

                Next, assume $4 \mid (5^{n-1}-1)$ for some integer $n > 1$.  By the definition of divides, we can write $$5^{n-1}-1=4k$$ for some integer $k$. Adding one to each side of the equation we get $$5^{n-1}=4k+1.$$
                Then $$5^n-1=5\cdot5^{n-1}-1 = 5(4k+1)-1 = 20k + 5 -1= 20k+4 = 4(5k+1)$$
                and so $4 \mid (5^n-1).$
            \end{proof}
        \item In the proof you copied, $P(n)$ was the statement $4 \mid (5^n-1)$.  What what the inductive hypothesis $P(n-1)$?  Hint: Look after the word ``assume''.
        \item Explain the algebra step-by-step in the last line of the proof.
    \end{actpart}
\end{act}

Let's look at more examples.

\begin{exam}[Number of bit strings]
\label{exam_math_ind_number_bs}
~
    \begin{apart}
        \item In \Cref{thm_number_bit_strings}, we proved that for $n \ge 0$ there are $2^n$ bit strings of length $n$.  Use mathematical induction to prove that there are $2^n$ bit strings of length $n$.
            \begin{soln}
                \emph{Proof} Induct on $n$. First, when $n=0$, there is exactly one bit string of length zero, namely $\lambda$, the empty bit string.  Since $2^0=1$, it follows that there are $2^0$ bit strings of length zero.

                Next, assume there are $2^{n-1}$ bit strings of length $n-1$ for some integer $n > 0$.  Let $b$ be a bit string of length $n$. Since the length of $b$ is at least 1, we can consider two cases.  
                
                Case 1: $b$ starts with \bs{0}.  Then $b=\bs{0}b_0$ where $b_0$ is a bit string of length $n-1$.  By the inductive hypothesis, there are $2^{n-1}$ choices for the bit string $b_0$.  Therefore, there are $2^{n-1}$ choices for the bit string $b$ in this case.  

                Case 2: $b$ starts with \bs{1}.  Then $b=\bs{1}b_1$ where $b_1$ is a bit string of length $n-1$.  By the inductive hypothesis, there are $2^{n-1}$ choices for the bit string $b_1$.  Therefore, there are also $2^{n-1}$ choices for the bit string $b$ in this case.

                Then, since cases add (\Cref{thm_cases_add}), there are $$2^{n-1}+2^{n-1}=2\cdot2^{n-1}=2^n$$ bit strings of length $n$. 
            \end{soln}
        \item What was the inductive hypothesis? 
            \begin{soln}
                The inductive hypothesis was that there are $2^{n-1}$ bit strings of length $n-1$ for some integer $n > 0$.
            \end{soln}
        \item Explain step-by-step why the algebra in the final equation is correct. 
            \begin{soln}
                For any number $x$ we have $x+x=x(1+1)=2x$.  In particular, $$2^{n-1}+2^{n-1}=2\cdot2^{n-1}.$$
                Next, 
                $2\cdot2^{n-1}= 2 \cdot \underbrace{2 \times 2 \times \cdots \times 2}_{(n-1) \text{ times}}=  \underbrace{2 \times 2 \times \cdots \times 2}_{n \text{ times}}=2^n.$
            \end{soln}
    \end{apart}
\end{exam}

We conclude with an example of a proof by mathematical induction about trees.  We use the fact that any tree with at least two vertices contains at least one leaf (degree one vertex).

\begin{exam}[Tree with $n$ vertices has exactly $n-1$ edges]
\label{exam_math_ind_tree_size}
    Use mathematical induction to prove that a tree with $n$ vertices has exactly $n-1$ edges.
        \begin{soln}
            \emph{Proof} Induct on $n$.  
            
            First, when $n=1$ the only tree is a single vertex which has zero edges, and $1-1=0$.  
            
            Next, assume that any tree with $(n-1)$ vertices has exactly $(n-1)-1 = n-2$ edges for some integer $n \ge 2$.  Let $T$ be a tree with $n$ vertices. Since every tree with at least two vertices has at least one leaf, there exists a vertex $v$ in $T$ such that $\fn{deg}(v)=1$.  
            
            Consider the graph $T'$ formed by removing $v$ and the one edge ending at $v$ from $T$.  Since $T$ was a tree, it does not contain any cycles so neither does $T'$.  Since $T$ is connected, any two vertices in $T$ are connected by a path in $T$.  Since $\fn{deg}(v)=1$, there is no path that includes $v$ in the middle (because that would mean $\fn{deg}(v)\ge2$).  Therefore any two vertices other than $v$ are connected by a path in $T$ that does not include $v$ and so $T'$ is still connected.  Last, $T'$ has $(n-1)$ vertices because we removed one vertex.  
            
            By our inductive hypothesis, it follows that $T'$ has exactly $(n-2)$ edges.  Since we only removed one edge from $T$, it follows that $T$ has exactly $(n-2)+1=n-1$ edges. 
        \end{soln}
\end{exam}

%%%%%%%%%%%%%%%%%%%%%%%%

\newcommand{\subalgind}{Algebra: Evaluating $f(n-1)$}
\subsection{\subalgind}
\label{sub_alg_ind}

Proofs by mathematical induction can involve a lot of algebra.  It would be a good moment to review:
    \begin{bull}
        \item \Cref{sub_integer_algebra} {\subintalgebra},
        \item \Cref{sub_factorials} {\subfactorials},
        \item \Cref{sub_exp_log} {\subexplog},
        \item \Cref{sub_fractions} {\subfractions}, and         
        \item \Cref{sub_polys} {\subpolys}.
    \end{bull}

Specifically, within proofs of mathematical induction, we often want to evaluate and simplify $f(n-1)$ where $f$ is a function and $n$ is an integer.  Let's see a few examples.

\begin{exam}[Evaluating $f(n-1)$]
\label{exam_eval_fofnminus1}
    In each part, evaluate $f(n-1)$ and simplify your answer.
        \begin{apart}            
            \item Consider the function $f: \mathbb{Z} \to \mathbb{Z}$ defined by $f(n)=2^{n+1}-1$.
                \begin{soln}
                    It may be helpful to think of $f(\Box) = 2^{(\Box)+1} -1$. To calculate $f(n-1)$, we replace $n$ in the original formula (or $\Box$ in this new version) with $(n-1)$ in parentheses to get
                    $$f(n-1) = 2^{(n-1)+1}-1=2^n-1.$$
                \end{soln}
            \item Consider the function $f: \mathbb{Z} \to \mathbb{Z}$ defined by $f(n)=n^2-n$.
                \begin{soln}
                    It may be helpful to think of $f(\Box) = (\Box)^2-(\Box).$ To calculate $f(n-1)$, we replace each occurrence of $n$ in the original formula (or $\Box$ in this new format) by $(n-1)$ in parentheses to get
                    $$f(n-1) = (n-1)^2-(n-1) = n^2-2n+1-n+1=n^2-3n+2$$
                \end{soln}
        \end{apart}
\end{exam}

Your turn to practice.

\begin{act}[Evaluating $f(n-1)$]
\label{act_eval_fofnminus1}  
    In each part, evaluate $f(n-1)$ and simplify your answer.
        \begin{actpart}
            \item $f: \mathbb{Z} \to \mathbb{Z}$ defined by $f(n)=3n+1$.
            \item $f: \mathbb{Z} \to \mathbb{Z}$ defined by $f(n)=n(n+1)$.
            \item $f: \mathbb{Z} \to \mathbb{Z}$ defined by $f(n)=3 \cdot 10^n$.
            \item $f: \mathbb{Z} \to \mathbb{Z}$ defined by $\displaystyle f(n)=a \left(\frac{r^{n+1}-1}{r-1}\right)$, for constants $a$ and $r$.
        \end{actpart}
\end{act}

%%%%%%%%%%%%%%%%%%%%%%%%

\newcommand{\subpffindseq}{Proof Format: Special Case of Mathematical Induction for Proving an Explicit Rule for a Recursively-defined Sequence}
\subsection{\subpffindseq}
\label{sub_pff_verify_explicit_seq}

One special case where proof by mathematical induction is used is to prove an explicit rule for a recursively-defined sequence.  We state this special case as its own proof format.

\begin{pff}[Mathematical induction: special case proving an explicit rule for a recursively-defined sequence]
\label{pff_math_ind_seq}
    Given the sequence $a$ defined recursively by $a_s=$ (initial condition) and $a_n=$ (recursive rule) for $n \ge s+1$, prove $a_n=$ (explicit rule) for $n \ge s$.
    \begin{proof}
        Write $f(n) = $(explicit rule).  Induct on $n$.

        First, when $n=s$ we know that $a_s=$ (initial condition) and also $f(s)=$ (evaluate and simplify to get the same answer).

        Next, assume $a_{n-1}=f(n-1)= $ (evaluate and simplify), for some integer $n > s$.  Then 
            \begin{align*}
                a_n & = \text{(recursive rule)} \\
                & = \cdots \text{(use the inductive hypothesis to replace } a_{n-1} \text{ by } f(n-1)) \cdots \\ 
                & = \cdots \text{(use algebra)} \cdots \\
                & = f(n).
            \end{align*}
    \end{proof}
\end{pff}

Note that in the base case we prove $a_s=f(s)$.  The inductive hypothesis is $a_{n-1}=f(n-1)$ for some integer $n > s$.  In the inductive step, we prove $a_n=f(n)$. 
As usual, here is an example of the proof format for you to copy and edit.

\begin{act}[Use mathematical induction to prove the explicit form of a recursively-defined sequence]
\label{act_pff_math_ind_seq}
~
    \begin{actpart}
        \item Copy the following proof.
        
        Consider the sequence $a$  defined recursively by $a_0= 1$ and $a_n = 2a_{n-1}$ for $ n \ge 1$.  Prove that $a_n=2^n$ for $n \ge 0$.
        
            \emph{Proof}  Write $f(n) = 2^n$. Induct on $n$.  
                
            First, when $n=0$ we know that $a_0 = 1$ and also $f(0) = 2^0 = 1$.  
                
            Next, assume $a_{n-1} = f(n-1)= 2^{n-1}$ for some integer $n \ge 0$.  Then 
                
            \hfill $a_n = 2a_{n-1} = 2 \cdot 2^{n-1} = 2^n = f(n)$ \hfill $\Box$
           
        \item Edit the proof to prove if the sequence $b$ is defined recursively by $b_0=10$ and $b_n = 3b_{n-1}$ for $n \ge 1$, then $b_n=10\cdot3^n$ for $n \ge 0$.
        \item The sequence $u$ is defined by $u_0= 10$ and $u_n = u_{n-1}^2$ for $n \ge 1$. Use mathematical induction to prove $u_n = 10^{2^n}$ for $n \ge 0$.  Be careful: $10^{2^n}$ means $10^{(2^n)}$.
    \end{actpart}
\end{act}

%%%%%%%%%%%%%%%%%%%%%%%%

\subsection{Exercises}

%%%%%%%%%%%%%%%%%%%%%%%%

\subsubsection*{Exercises for \Cref{sub_pff_math_induction}
 \subpffmathind}

\begin{exer}[\exerP] % Divisibility proof by induction
\label{exer_divisibility_proof_induction}
    Use mathematical induction to prove that $3 \mid (4^n-1)$ for all integers $n \ge 1$.  
\end{exer}

\begin{exer}[\exerU] % prove every integer is even or odd]
\label{exer_math_ind_every_int_even_or_odd}
    Use mathematical induction to prove every integer is either even or odd.  Use that $n=(n-1)+1$ and the inductive hypothesis to break into two cases: $n-1$ is even or $n-1$ is odd.
\end{exer}

\begin{exer}[\exerU] % Counting edges in a complete graph
\label{exer_count_edges_Kn_proof_ind}
    Consider the function $f: \mathbb{Z} \to \mathbb{Z}$ defined by $f(n) = \frac{n(n-1)}{2}$.
        \begin{exerpart}
            \item Evaluate and simplify $f(1)$.
            \item Evaluate and simplify $f(n-1)$. Leave the numerator in factored form.
            \item  Copy the following outline of a proof by mathematical induction that the complete graph $K_n$ has $f(n)$ edges for $n \ge 1$ and fill in the missing details.

            \begin{proof}
                First, when $n=1$, the graph $K_1$ is one vertex which has 0 edges and $f(1)=\underline{\hspace{1in}}$.  Thus, $K_1$ has $f(1)$ edges.

                Next, assume $\underline{\hspace{1in}}$ for some integer $n > 1$.  The graph $K_n$ can be built from $K_{n-1}$ by adding one more vertex, $v$, and adding an edge from $v$ to each of the vertices of $K_{n-1}$.  Notice that we are adding $n-1$ new edges.  Thus, the number of edges in $K_n$ equals the number of edges in $K_{n-1}$ and $n-1$ new edges. That is, the number of edges of $K_n$ is $$f(n-1) + (n-1) = \underline{\hspace{.25in} \text{ (show algebra) }\hspace{.25in}} = f(n).$$
            \end{proof}
    \end{exerpart}
\end{exer}

\begin{exer}[\exerU] %Number of subsets of a set.
\label{exer_math_ind_number_subsets}
    Use mathematical induction to prove that there are $2^n$ subsets of the set $\{1,2,3,\ldots,n\}$.  Hint: Look at \Cref{exam_math_ind_number_bs}.  In the inductive step, break the count into two cases based on whether the subset contains $n$. If the subset $S$ does not contain $n$, then $S \subseteq \{1,2,3, \ldots, n-1\}$.  If the subset $S$ contains $n$, then remove $n$ from the subset to get $S - \{n\} \subseteq \{1,2,3, \ldots, n-1\}$.
\end{exer}

\begin{exer}[\exerR] % proof format mathematical induction
\label{exer_dyk_pff_math_induction}
    Do you know \ldots
        \begin{exerpart}
            \item When we might use a proof by mathematical induction?
            \item Why proof by mathematical induction works (logically)?
            \item Why we need the base case?
            \item What we prove in the base case of a proof by mathematical induction?
            \item What the inductive hypothesis is in a proof by mathematical induction?
            \item What we prove in the inductive step of a proof by mathematical induction?
        \end{exerpart}
\end{exer}

 \begin{exer}[\exerE] % number of permutations of $\{1,2,3,\ldots,n\}$]
\label{exer_math_ind_number_perms}
    Use mathematical induction to prove that there are $n!$ permutations of the set $\{1,2,3,\ldots,n\}$.  Recall that a \vocab{permutation} of the set $\{1,2,3,\ldots,n\}$ is an $n$-tuple that uses each of $1,2,3,\ldots,n$ exactly once.  Hint: Given a permutation $p$ of $\{1,2,3,\ldots,n\}$ we can remove $n$ to get a permutation $p_0$ of $\{1,2,3,\ldots,n-1\}$.  There are $(n-1)!$ such resulting permutations and there are $n$ choices for the position where $n$ appeared in $p$: directly before the 1, directly before the 2, directly before the 3, \ldots, directly before the $n-1$, or at the very end.  Use that steps multiply (\Cref{thm_steps_multiply}).
\end{exer}

%%%%%%%%%%%%%%%%%%%%%%%%

\subsubsection*{Exercises for \Cref{sub_alg_ind} \subalgind}
%   %%%%%%%%%%%%%%%%%%
% \subsubsection{\answers for \Cref{sub_alg_ind} \subalgind}

 \begin{exer}[\exerP] % evaluating $f(n-1)$ involving polynomials]
\label{exer_eval_fofnminus1_linears}
    In each part, evaluate $f(n-1)$ and simplify your answer.
        \begin{exerpart}
            \item $f: \mathbb{Z} \to \mathbb{Z}$ defined by $\displaystyle f(n)=\frac{2n}{n+1}$.
            \item $f: \mathbb{Z} \to \mathbb{Z}$ defined by $\displaystyle f(n)=\frac{n}{2n+1}$.
            \item $f: \mathbb{Z} \to \mathbb{Z}$ defined by $\displaystyle f(n)=\frac{n(n+1)}{2}$. Leave the answer in factored form.
        \end{exerpart}
\end{exer}

\begin{exer}[\exerP] % evaluating $f(n-1)$ involving exponentials]
\label{exer_eval_fofnminus1_exponentials}
    In each part, evaluate $f(n-1)$ and simplify your answer.
        \begin{exerpart}
            \item $f: \mathbb{Z} \to \mathbb{Z}$ defined by $f(n)=(-1)^n$.
            \item $f: \mathbb{Z} \to \mathbb{Z}$ defined by $f(n)=2\cdot3^{n+1}-1$.
            \item $f: \mathbb{Z} \to \mathbb{Z}$ defined by $f(n) =10^{2^n}$
        \end{exerpart}
\end{exer}

\begin{exer}[\exerP] % Evaluating $f(n-1)$ for a variety of functions
\label{exer_eval_fofnminus1_exponentials}
    In each part, evaluate $f(n-1)$ and simplify your answer.
        \begin{exerpart}
            \item $f: \mathbb{Z} \to \mathbb{Z}$ defined by $f(n)=n^2+n+1$.
            \item $f: \mathbb{Z} \to \mathbb{Z}$ defined by $f(n)=n(n+1)(2n+1)$.  Leave your answer factored.
            \item $f: \mathbb{Z} \to \mathbb{Z}$ defined by $f(n)=(n+1)!-1$.
        \end{exerpart}
\end{exer}

\begin{exer}[\exerU] % simplifying algebraic expressions
\label{exer_simplify_ind_alg}
    Simplify each expression.
        \begin{exerpart}
            \item $(2(n-1)+1)+2$
            \item $2\cdot (2^n-1)+1$
            \item $n(n-1)!$
        \end{exerpart} 
\end{exer}

\begin{exer}[\exerU] % simplifying algebraic expressions with fraction
\label{exer_simplify_ind_alg_frac}
    Simplify each expression.  To add the fractions, write over the common denominator 2.
        \begin{exerpart}
           \item $\displaystyle 3\left( \frac{3^n-1}{2} \right) + 1$
            \item $\displaystyle \frac{(n-1)n}{2} +n$
        \end{exerpart} 
\end{exer}

\begin{exer}[\exerR] % evaluating f(n-1)
\label{exer_dyk_fofnminus1_induction}
    Do you know \ldots
        \begin{exerpart}
            \item How to evaluate $f(n-1)$ for a function $f$ and why we might need parentheses?
            \item How to simplify $2\cdot2^{n-1}$ and other related exponential expressions?
            \item How to simplify $n(n-1)!$ and other related expressions involving factorials? 
        \end{exerpart}
\end{exer}

%%%%%%%%%%%%%%%%%%%%%%%%

\subsubsection*{Exercises for \Cref{sub_pff_verify_explicit_seq}
 \subpffindseq}

\begin{exer}[\exerP] % using mathematical induction to prove explicit rule for geometric sequence]
\label{exer_math_ind_geom}
    The sequence $a$ is defined by $a_0= 3$ and $a_n = 10a_{n-1}$ for $n \ge 1$. Use mathematical induction to prove $a_n = 3\cdot 10^n$ for $n \ge 0$.
\end{exer}

\begin{exer}[\exerP] % using mathematical induction to prove explicit rule for arithmetic sequence]
\label{exer_math_ind_arith}
    The sequence $b$ is defined by  $b_0= 1$ and $b_n = b_{n-1} + 3$ for $n \ge 1$. Use mathematical induction to prove $b_n = 3n+1$ for $n \ge 0$.
\end{exer}

\begin{exer}[\exerP] % using mathematical induction to prove explicit rule for sequence involving alternating signs]
\label{exer_math_ind_alt}
    The sequence $t$ is defined by $t_0= 5$ and $t_k = -t_{k-1}$ for $k \ge 1$. Use mathematical induction to prove $t_k = (-1)^k5$ for $k \ge 0$.
\end{exer}

\begin{exer}[\exerP] % using mathematical induction to prove explicit rule for sequence involving exponential]
\label{exer_math_ind_exp}
    The sequence $c$ is defined by $c_0=1$ and $c_k = 2c_{k-1}+3$ for $n \ge 1$. Use mathematical induction to prove $c_k = 2^{k+2}-3$ for $k \ge 0$.  We saw this sequence in \Cref{exam_double_and_add3}.
\end{exer}

\begin{exer}[\exerU] % using mathematical induction to prove explicit rule for sequence involving fraction]
\label{exer_math_ind_frac}
    The sequence $s$ is defined by $s_0=0$ and $s_m = 5s_{m-1}+1$ for $m \ge 1$. Use mathematical induction to prove $s_m = \frac{5^m-1}{4}$ for $m \ge 0$
\end{exer}

\begin{exer}[\exerU] % using mathematical induction to prove explicit rule for handshakes sequence]
\label{exer_math_ind_hs}
    The sequence $u_n$ is defined by $u_1= 1$ and $u_n = u_{n-1}+n$ for $n \ge 2$. Use mathematical induction to prove $u_n =  \frac{n(n+1)}{2}$ for $n \ge 1$.
\end{exer}

\begin{exer}[\exerU] % using mathematical induction to prove explicit rule for factorial sequence]
\label{exer_math_ind_factorial}
    The sequence $v_n$ is defined by $v_1= 1$ and $v_n = nv_{n-1}$ for $n \ge 2$. Use mathematical induction to prove $v_n = n!$ for $n \ge 1$.
\end{exer}

\begin{exer}[\exerU] % using mathematical induction to prove explicit rule for quadratic sequence]
\label{exer_math_ind_quadratic}
    The sequence $x_n$ is defined by $x_0=1$ and $x_n = x_{n-1}+2n$ for $n\ge 1$. Use mathematical induction to prove $x_n =n^2 + n + 1$ for $n \ge 0$
\end{exer}

\begin{exer}[\exerR] % verifying explicit formulas for sequences using mathematical induction
\label{exer_dyk_verify_explicit_induction}
    Do you know \ldots
        \begin{exerpart}
            \item How to use mathematical induction to prove an explicit rule for a recursively-defined sequence?
            \item What the base case is in such proofs?
            \item What the inductive hypothesis is in such proofs?
            \item Which is used first in the final line -- the recursion or the inductive hypothesis?
        \end{exerpart}
\end{exer}

\begin{exer}[\exerE] % savings account balance formula]
\label{exer_savings_balance}
    We deposit \$\text{10,000} into a bank account that earns 3\% interest annually. Write $a_n=$ the account balance at the end of $n$ years. We saw in \Cref{exer_savings_annuity} (\ref{part_savings}) that $a_0=\text{10,000}$ and $a_n=1.03a_{n-1}$ for $n \ge 1$. Let's generalize.  If we deposit \$$p$ into a bank account that earns $r$ interest annually, where $r$ is the percentage written as a decimal, then $a_0=p$ and $a_n=(1+r)a_{n-1}$ for $n \ge 1$.  Use mathematical induction to prove $a_n=p(1+r)^n$ for $n \ge 0$.
\end{exer}

\begin{exer}[\exerE] % annuity balance formula]
\label{exer_annuity_balance}
    We deposit \$\text{10,000} into a bank account that earns 3\% interest annually but then we deposit another \$\text{10,000} at the end of each year (just after interest is paid).  Write $b_n=$ the account balance at the end of $n$ years.   We saw in \Cref{exer_savings_annuity} (\ref{part_annuity}) that $b_0=\text{10,000}$ and $b_n = 1.03b_{n-1}+ \text{10,000}$.  Let's generalize.  If we deposit \$$p$ into a bank account that earns $r$ interest annually, where $r$ is the percentage written as a decimal, but then we deposit another \$$p$ at the end of each year (just after interest is paid), then $b_0=p$ and $b_n = (1+r)b_{n-1}+ p$. Use mathematical induction to prove that $b_n = p \displaystyle \left(\frac{(1+r)^{n+1}-1}{r}\right)$ for $n \ge 0$.
\end{exer}

%%%%%%%%%%%%%%%%%%%%%%%%%%%
\subsection{\answers}
%%%%%%%%%%%%%%%%%%%%%%%%%%%
 %%%%%%%%%%%%%%%%%%
\subsubsection{\answers for \Cref{sub_pff_math_induction}
 \subpffmathind}

\subsubsection*{\Cref{exer_divisibility_proof_induction} \answers}
Hint: Copy and edit the proof from \Cref{act_div_proof_math_ind}.

\subsubsection*{\Cref{exer_count_edges_Kn_proof_ind} \answers}
\begin{apart}
    \item \label{part_edges_K1_math_ind_ans} Hint: $1-1=0$
    \item \label{part_edges_Kn_math_ind_ans} $\frac{(n-1)(n-2)}{2}$
    \item Hint for the first blank: Fill in the answer to (\ref{part_edges_K1_math_ind_ans}).  
    
    Hint for the second blank: $P(n)$ is ``The graph $K_n$ has $f(n)$ edges.''  Fill in the blank with $P(n-1)$.  
    
    Hint for the third blank:  First, use your answer to (\ref{part_edges_Kn_math_ind_ans}), then write $n-1 = \frac{2(n-1)}{2}$ so the fractions have a common denominator, and then add the fractions.
\end{apart}

\subsubsection*{\Cref{exer_eval_fofnminus1_linears} \answers}
The answers, in some order, are
$\frac{n-1}{2n-1}$, $\frac{n(n-1)}{2}$, and $\frac{2n-2}{n}$.

\subsubsection*{\Cref{exer_eval_fofnminus1_exponentials} \answers}
\begin{apart}
    \item Hint: If you evaluate your answer at $n=3$, you should get 7. If you did not get that, check that you calculated $(n-1)^2=(n-1)(n-1)$ correctly.
    \item Hint: One factor is $2n-1$.  If you did not get that, check that you entered $(n-1)$ in parentheses.
    \item Hint: It does not simplify much.
\end{apart}

\subsubsection*{\Cref{exer_simplify_ind_alg_frac} \answers}
\begin{apart}
    \item Hint: Your answer should be one fraction with denominator 2.  It should involve $3^{n+1}$.  
    \item Hint: $n = \frac{2n}{2}$
\end{apart}

   %%%%%%%%%%%%%%%%%%
\subsubsection{\answers for \Cref{sub_pff_verify_explicit_seq}
 \subpffindseq}

\subsubsection*{\Cref{exer_math_ind_geom} \answers}
Hint: Copy and edit the proof from \Cref{act_pff_math_ind_seq}.

\subsubsection*{\Cref{exer_math_ind_arith} \answers}
Hint: Copy and edit the proof from \Cref{act_pff_math_ind_seq}.  When you evaluate $f(n-1)$, make sure $(n-1)$ is in parentheses.

\subsubsection*{\Cref{exer_math_ind_exp} \answers}
Hint: Notice that we are using $k$ instead of $n$.  Some useful algebra facts are $2\cdot 2^{k+1}=2^{k+2}$ and $2(-3)=-6$.

\subsubsection*{\Cref{exer_math_ind_frac} \answers}
Hint: Notice that we are using $m$ instead of $n$. Some useful algebra facts are $5(5^m-1) = 5^{m+1}-5$ and $1 = \frac{4}{4}$.

\subsubsection*{\Cref{exer_math_ind_hs} \answers}
Hint: Some useful algebraic facts are $(n-1)n = n^2-n$, $n = \frac{2n}{2}$, and $n^2+n=n(n+1)$.

\subsubsection*{\Cref{exer_math_ind_factorial} \answers}
Hint: $7\cdot 6! = (7)(6)(5)(4)(3)(2)(1)=7!$  Use this idea to simplify $n(n-1)!$

\subsubsection*{\Cref{exer_math_ind_quadratic} \answers}
Hint: Make sure you have correctly simplified $(n-1)^2=(n-1)(n-1)$.


